\let\oldprintchaptertitle=\printchaptertitle
\renewcommand{\printchaptertitle}[1]{%
	\vspace*{-75pt}
	\oldprintchaptertitle{#1}
}%
\myChapterStar{Résumé}{}{section}
\let\printchaptertitle=\oldprintchaptertitle
Dans un contexte technologique en perpétuelle évolution, les architectures orientées services (SOA) et les microservices se sont imposés comme des paradigmes incontournables pour concevoir des systèmes logiciels modulaires, flexibles et distribués. 
Ces architectures permettent de décomposer les fonctionnalités d'une application en services autonomes, facilitant ainsi la réutilisabilité, l'évolutivité et l’adaptabilité. Cependant, la manière dont ces services sont composés reste un enjeu central pour garantir performance, robustesse et réactivité, surtout dans les environnements dynamiques tels que le cloud computing ou les systèmes orientés données.
Les approches classiques de composition – \textit{statique ou dynamique} – montrent leurs limites face aux contraintes modernes de scalabilité, de résilience et de traitement paresseux. 
% C’est dans ce cadre que s’inscrit la composition incrémentale des services, une approche émergente orientée données, où l'exécution d’un processus s'effectue de manière progressive, guidée par la disponibilité des informations locales. Elle s’appuie notamment sur le modèle des  Grammaires Attribuées Gardées (GAG), qui permet de représenter la logique d’assemblage de services sous forme de règles pilotées par des dépendances explicites entre attributs.
C’est dans ce cadre que s’inscrit la composition incrémentale des services, une approche émergente orientée données, où l’exécution d’un processus s’effectue de manière progressive, guidée par la disponibilité des informations calculées et attendues par les services du processus. Elle s’appuie notamment sur le modèle des Grammaires Attribuées Gardées (GAG), qui permet de représenter la logique de composition de services sous forme de règles pilotées par des dépendances explicites entre les données. La particularité de la composition incrémentale de services est qu’elle supporte la conteneurisation, qui est largement utilisée dans les architectures cloud et par les entreprises mainstream telles que Google, Amazon et IBM. Ceci la démarque des approches traditionnelles basées sur des modèles ad hoc ne supportant pas la conteneurisation.
%----->
% Afin d’évaluer rigoureusement cette approche, nous avons propose un benchmark dédié à la transformée de Fourier incrémentale, déployé sur un cluster Kubernetes d'une machine physique. Les résultats expérimentaux montrent que la composition incrémentale offre une meilleure adaptabilité à l’environnement d’exécution, un parallélisme naturel mieux exploité, et une robustesse accrue face aux défaillances partielles, ce qui en fait une approche prometteuse pour les systèmes distribués modernes.
Afin d’évaluer rigoureusement cette approche, nous avons proposé un benchmark dédié à la transformée de Fourier incrémentale, déployé sur un cluster Kubernetes d’une machine physique.
%---->
% Cette étude ouvre la voie à de nouvelles pistes pour l’optimisation des systèmes à base de services paresseux, et souligne l’intérêt d’un couplage étroit entre modélisation formelle (GAG) et infrastructures modernes (Kubernetes).
Les résultats expérimentaux montrent que la composition incrémentale offre une meilleure adaptabilité à l’environnement d’exécution et un parallélisme naturel qui réduit significativement le temps d’exécution. Bien plus, elle permet de proposer des implémentations plus robustes qui traitent de plus grandes quantités de données tout en offrant une meilleure résilience aux pannes. Nos résultats se généralisent à tout autre algorithme suivant le schéma DPR (Diviser pour régner).
% Ce travail met ainsi en lumière l’intérêt de combiner des modèles formels comme les GAG avec des technologies cloud-native telles que Docker, Kubernetes et Vagrant, pour concevoir des moteurs de composition efficaces, intelligents et adaptatifs.

\vspace{1cm}
\noindent\textbf{Mots clés:} Architecture Orientée services, Grammaires Attribuées Gardées, Services paresseux, composition incrementale, benchmark, transformée de Fourier, Kubernetes, cloud computing.

\myCleanStarChapterEnd

